\begin{center}
{\Large \textbf{Abstract}}
\end{center}
Recent progress in natural language processing has been largely driven by large language models(LLMs).
These architectures aim to learn representations of text that capture syntax, 
semantics, and long-range dependencies in an efficient and scalable way. 
Despite significant advances in performance, the internal mechanisms by which these models store and manipulate information 
remain only partially understood. \\ \\
% 
This limitation becomes especially relevant in deep models that rely on fundamentally 
different approaches to information storage. Transformer-based models process sequences using 
self-attention, where information is mixed across tokens at each layer without an explicit notion 
of persistent memory. In contrast, the xLSTM, a recently developed architecture was proposed 
by Beck et al. \cite{beck2024xlstmextendedlongshortterm} as an alternative to the transformer architectures 
by extending the classic LSTM cell and reviving the interest in Recurrent Neural Networks \\ \\ 
%
This work performs an analysis of the xLSTM architecture through a set of interpretability 
experiments aimed at understanding how information is represented, stored, and transformed inside the 
model. Moreover, the practicality of this architecture is assessed by evaluating it in a question answering 
setup, a very common application of LLMs. \\ \\ 
%
We track how the model's intermediate hidden states evolve and compare them to those produced by a Transformer model
to assess the similarity of their internal representations. 
Furthermore, we inspect memory behavior during the processing of sequences to understand how individual tokens influence 
the retention and discarding information. \\ \\
%
Finally, we evaluate memory caching as a practical use of xLSTM's fixed-size memory state. By reusing the memory 
state after processing the context, the model avoids recomputing the input sequence and achieves inference 
speedups that increase with context length. We compare this approach against Mistral-7B and examine the 
trade-off between inference efficiency and task performance.


